\section{Related Work}

\subsection{Alignment and Preference Optimization}
Optimizing Large Language Models (LLMs) to align with human intent has predominantly relied on Reinforcement Learning from Human Feedback (RLHF) \cite{ouyang2022training}, utilizing Proximal Policy Optimization (PPO) \cite{schulman2017proximal}. Recently, Direct Preference Optimization (DPO) \cite{rafailov2023direct} emerged as a stable alternative that mathematically bypasses the need for an explicit reward model. Building on this, recent advancements like RRHF \cite{yuan2023rrhf}, ORPO \cite{hong2024orpo}, and SimPO \cite{meng2024simpo} have further simplified alignment via reference-free objectives. The state-of-the-art in 2025 has moved toward adaptive batch-wise scheduling \cite{zhang2025adaptive} and iterative dual-based methods for constrained alignment \cite{chen2025constrained}, addressing the inherent algorithmic biases identified in standard RLHF \cite{ji2024limits, shrivastava2024algorithmic}.

\subsection{Hallucination Mitigation and Reasoning}
In knowledge-intensive tasks, such as educational explanation generation, factual correctness is paramount. Prior work has explored grounding LLMs using retrieval-augmented generation (RAG) \cite{lewis2020retrieval} or process-based feedback \cite{uesato2023verify}. Modern techniques focus on mechanistic hallucination detection \cite{li2023haludetect} and Multi-Model Contrastive Decoding (MCD) \cite{liu2025contrastive}. While debates continue regarding whether emergent abilities are a metric-driven "mirage" \cite{schaeffer2023mirage}, advancements in 2025 demonstrate that reasoning quality can be significantly enhanced by verifying "decision pivots" in Chain-of-Thought traces \cite{wang2025pivots} or equipping models with cognitive tools \cite{zhao2025cognitive}. However, these logical optimizations often degrade fluency—the "alignment tax"—which our work addresses via a hybrid verifier-consistent reward.

\subsection{LLM-as-a-Judge and Evaluation Bias}
The use of powerful instruction-tuned models like GPT-4 as automated evaluators \cite{zheng2024judging} is widespread but prone to significant biases, including verbosity bias \cite{wang2023large} and self-preference bias \cite{pan2024evalbias}. Comprehensive assessments like DecodingTrust \cite{wang2023trust} highlight the vulnerability of current models to stylistic manipulation. Our Hybrid-DPO framework circumvents these issues by fusing objective NLI metrics with constrained fluency verifiers, effectively balancing the trade-off documented in recent $\beta$-DPO variants \cite{park2024betadpo} while maintaining safety and knowledge retention standards \cite{kim2025safety}.
